\begin{tikzpicture}[x={\dimexpr\linewidth/169\relax},y=1mm,
  font=\figurefont\footnotesize,text=NNInk]
  \nnTensorState{x}{9}{0}{NNInput}{$\bm x$}
  \nnConvOp{f}{42}{0}{3\times3}
  \nnReluOp{act}{70}{0}
  \nnConvOp{g}{42}{-28}{3\times3}
  \nnSigmoidOp{sig}{70}{-28}
  \node[nn pointwise] (mul1) at (105,0) {$\odot$};
  \node[nn operator,minimum width=12mm,minimum height=10mm]
    (comp) at (105,-28) {$1-\,\cdot$};
  \node[nn pointwise] (mul2) at (119,-53) {$\odot$};
  \node[nn pointwise] (add) at (137,0) {$+$};
  \nnTensorState{y}{160}{0}{NNOutput}{$\bm y$}
  \draw[nn flow] (x.east)--(f.west);
  \draw[nn operator path] (f.east)--(act.west);
  \draw[nn operator path] (act.east)--node[above=1.5mm] {$\bm F(\bm x)$} (mul1.west);
  \draw[nn operator path] (mul1.east)--(add.west);
  \draw[nn operator path,draw=NNOutput] (add.east)--(y.west);
  % 共享输入主干连续穿过两个分支点，只在底部转弯处使用圆角。
  \draw[nn operator path,draw=NNInput,line width=1.1pt]
    (24,0)--(24,-53)--(mul2.west);
  \draw[nn flow] (24,-28)--(g.west);
  \draw[nn operator path] (g.east)--(sig.west);
  \draw[nn operator path] (sig.east)--node[below=1.5mm] {$\bm g$} (comp.west);
  \draw[nn operator path] (87,-28)--(87,-14)--(105,-14)--(mul1.south);
  \draw[nn operator path] (comp.east)--(119,-28)--(mul2.north);
  \draw[nn operator path] (mul2.east)--(137,-53)--(add.south);
  \foreach \yy in {0,-28}{\fill[NNInput] (24,\yy) circle[radius=.5mm];}
  \fill[NNHidden] (87,-28) circle[radius=.5mm];
\end{tikzpicture}
